\section{Scope and open questions}

The product theorems give a complete answer to two construction questions:
the witness pattern combines by Boolean OR, and group isotopy holds exactly
when it holds in both factors. They do not classify all FFF squares. The
explicit nongroup loop shows that group tables are not the only source,
and the order-$12$ witness rules out a power-of-two-only existence spectrum.

The existence of all eight patterns in orders $2^m$, $m\geq3$, uses finite
order-$8$ starting examples. The nongroup FFF family in the same orders
instead follows from the explicit loop and written proofs. No statement
here counts distinct higher-order main classes or assumes the product map
is injective on them.

Unrestricted order $18$ remains open. A failed construction, a timeout, or
an exclusion for one parametrized family is not a nonexistence proof.
The full report's separately evidenced order-$10$ exclusion is not repeated
in this paper. Further families, represented-class lower bounds, and
bounded order-$18$ exclusions belong in the longer editions, where their
hypotheses and computational dependencies can be stated in full.

The next mathematical questions are which additional constructions yield
FFF squares, how restrictive the three views are together, and which
existence orders can be decided rigorously. Before any journal submission,
independent expert review and literature comparison remain necessary to
assess correctness, significance, and priority. This short selection is
intended to make that review manageable, not to substitute for it.
